\section*{Розділ VI. Документообіг та оприлюднення}
\addcontentsline{toc}{section}{Розділ VI. Документообіг та оприлюднення}

\subsection*{6.1. Ведення протоколів засідань}
\addcontentsline{toc}{subsection}{6.1. Ведення протоколів засідань}
    6.1.1. Хід кожного засідання КСУ фіксується у протоколі, який є офіційним документом, що підтверджує факт проведення засідання та прийняті рішення.

    6.1.2. Протокол засідання КСУ ведеться постійним секретарем КСУ, який призначається КСУ на строк повноважень складу делегатів.

    6.1.3. Протокол засідання КСУ містить:

        \begin{enumerate}[label=\alph*)]
            \item Дату, час та місце (форму) проведення засідання;
            \item Список присутніх делегатів КСУ та запрошених осіб;
            \item Затверджений порядок денний;
            \item Основні тези доповідей з кожного питання порядку денного;
            \item Стислий виклад ходу обговорення кожного питання;
            \item Детальні результати голосування (кількість голосів "за", "проти", "утрималися");
            \item Повний текст прийнятих рішень;
            \item Інші суттєві обставини засідання.
        \end{enumerate}

    6.1.4. Протокол засідання підписується електронним підписом Спікером КСУ та постійним секретарем КСУ не пізніше ніж через 3 робочі дні після проведення засідання.

    6.1.5. У разі відсутності постійного секретаря КСУ, його функції виконує секретар засідання, призначений Спікером КСУ з числа присутніх делегатів на початку засідання.

\subsection*{6.2. Оформлення рішень КСУ}
\addcontentsline{toc}{subsection}{6.2. Оформлення рішень КСУ}
    6.2.1. Рішення КСУ оформлюються як окремі документи або як невід'ємна частина протоколу засідання відповідно до вимог пункту 4.7 цього Регламенту.

    6.2.2. Кожне рішення КСУ має унікальний номер, що складається з порядкового номера засідання та порядкового номера рішення на цьому засіданні (наприклад: "Рішення КСУ №3/2024-5", де 3 -- номер засідання, 2024 -- рік, 5 -- номер рішення на засіданні).

    6.2.3. Відповідальність за правильне оформлення рішень КСУ покладається на Спікера КСУ та постійного секретаря КСУ.

    6.2.4. Всі рішення КСУ підписуються електронним підписом уповноважених осіб відповідно до вимог пункту 4.7.2 цього Регламенту.

    6.2.5. Тексти рішень КСУ перевіряються на відповідність обговоренню та результатам голосування перед їх підписанням та оприлюдненням.

\subsection*{6.3. Порядок оприлюднення рішень та протоколів}
\addcontentsline{toc}{subsection}{6.3. Порядок оприлюднення рішень та протоколів}
    6.3.1. Всі рішення КСУ та протоколи засідань підлягають обов'язковому оприлюдненню відповідно до принципу гласності діяльності КСУ.

    6.3.2. Строки оприлюднення документів КСУ:

        \begin{enumerate}[label=\alph*)]
            \item Протоколи засідань -- протягом 5 робочих днів з дня проведення засідання;
            \item Рішення КСУ -- протягом 3 робочих днів з дня їх прийняття;
            \item Невідкладні рішення КСУ -- протягом 1 робочого дня з дня їх прийняття.
        \end{enumerate}

    6.3.3. Оприлюднення здійснюється шляхом:

        \begin{enumerate}[label=\alph*)]
            \item Розміщення на офіційних інформаційних ресурсах органів студентського самоврядування Університету;
            \item Надсилання копій документів органам студентського самоврядування, які є виконавцями прийнятих рішень або до компетенції яких належать розглянуті питання;
            \item Інформування студентської спільноти через доступні канали зв'язку.
        \end{enumerate}

    6.3.4. Частини протоколу, що містять персональну інформацію або відомості з обмеженим доступом, можуть не підлягати повному оприлюдненню за рішенням Спікера КСУ.

\subsection*{6.4. Зберігання документів КСУ}
\addcontentsline{toc}{subsection}{6.4. Зберігання документів КСУ}
    6.4.1. Всі документи КСУ (протоколи засідань, рішення, супровідні матеріали) зберігаються в електронному архіві органів студентського самоврядування.

    6.4.2. Документи КСУ зберігаються протягом всього періоду діяльності відповідного складу КСУ та не менше 3 років після його завершення.

    6.4.3. Організаційне забезпечення функціонування електронного архіву покладається на Студентську раду КАІ.

\subsection*{6.5. Надання інформації та документів}
\addcontentsline{toc}{subsection}{6.5. Надання інформації та документів}
    6.5.1. Органи студентського самоврядування, їх посадові особи та студенти мають право звертатися до КСУ за отриманням копій документів КСУ, що стосуються їхніх прав та інтересів.

    6.5.2. Запити про надання документів КСУ подаються Спікеру КСУ у письмовій формі (включаючи електронну пошту або повідомлення у месенджерах) з зазначенням:

        \begin{enumerate}[label=\alph*)]
            \item Найменування запитуваних документів;
            \item Обґрунтування потреби в отриманні документів;
            \item Контактної інформації для надання відповіді.
        \end{enumerate}

    6.5.3. Спікер КСУ забезпечує надання запитуваних документів протягом 5 робочих днів з дня отримання запиту, якщо документи не містять інформації з обмеженим доступом.

    6.5.4. У наданні документів може бути відмовлено, якщо:

        \begin{enumerate}[label=\alph*)]
            \item Запитувані документи містять персональну інформацію третіх осіб;
            \item Документи стосуються питань, розглянутих на закритому засіданні КСУ;
            \item Надання документів може завдати шкоди інтересам студентського самоврядування або окремих осіб.
        \end{enumerate}

    6.5.5. Рішення про відмову у наданні документів приймається Спікером КСУ мотивовано та може бути оскаржене до КСУ.

\subsection*{6.6. Технічне забезпечення документообігу}
\addcontentsline{toc}{subsection}{6.6. Технічне забезпечення документообігу}
    6.6.1. Для забезпечення ефективного документообігу КСУ використовуються сучасні електронні засоби та платформи.

    6.6.2. Технічні вимоги до документообігу КСУ:

        \begin{enumerate}[label=\alph*)]
            \item Всі документи зберігаються у електронній формі з можливістю пошуку та фільтрації;
            \item Забезпечується резервне копіювання документів;
            \item Ведеться журнал доступу до документів з фіксацією дати, часу та особи, яка отримала доступ;
            \item Документи захищаються від несанкціонованого доступу та зміни.
        \end{enumerate}

    6.6.3. Відповідальність за технічне забезпечення документообігу КСУ покладається на Студентську раду КАІ у взаємодії зі Спікером КСУ.

    6.6.4. У разі технічних збоїв або неможливості використання електронних засобів, документообіг може тимчасово здійснюватися у паперовій формі з подальшим переведенням у електронний формат.